% !TeX root = ../main.tex

\section{Background and Related Work}\label{sec:background}

This section synthesises prior research to situate crypto key opinion leaders (KOLs) as a distinct category of digital workers. We organise the literature around four themes: (1) influencers, finfluencers, and creator labour; (2) why crypto constitutes a qualitatively different context; (3) crypto KOLs as digital creators under financialised conditions; and (4) the research gap motivating our study.

\subsection{Influencers, Finfluencers, and Creator Labour}

\subsubsection{Influencer Culture and Creator Labour}

Influencer research conceptualises social media creators as micro-celebrities who strategically curate authenticity and intimacy to attract audiences \cite{MarwickStatusUpdate2013,Khamis2017,Enke2019}. Ethnographic accounts reveal how influencers cultivate personas that are simultaneously relatable and aspirational, with authenticity understood as performative, carefully crafted to sustain attention and trust \cite{MarwickStatusUpdate2013,BaymPlaying2018}.

Two concepts capture the digital labour involved. Abidin's \emph{visibility labour} describes the ongoing work of posting, curation, and interaction through which influencers maintain presence \cite{AbidinVisibility2016}. Baym's \emph{relational labour} emphasises the effort invested in sustaining intimate audience relationships \cite{BaymPlaying2018}. Both are consequential because parasocial interaction makes followers receptive to recommendations; credibility emerges from affective bonds, not just expertise \cite{Lou2019,Labrecque2014,Reinikainen2020}.

From a political-economy perspective, Duffy and colleagues describe  ``nested precarities'' of creators, which denotes overlapping insecurities at the level of platforms, industries, and algorithms \cite{DuffyPlatformPractices2019,DuffyNestedPrecarities2021}. Creators engage in constant self-monitoring, cross-posting, and brand management to remain visible \cite{DuffyAspirationalWork2017}, ``playing the visibility game'' to satisfy opaque ranking systems \cite{CotterVisibilityGame2019}. Success depends on algorithmic regimes and revenue streams that can shift abruptly.

\subsubsection{Finfluencers and Social-Media Finance}

Recent scholarship extends influencer culture into personal finance. Guan characterises ``finfluencers'' as online financial advice-givers who mix analysis, storytelling, and entertainment, yet often operate outside regulatory frameworks \cite{GuanFinfluencer2023}. Hayes and Ben-Shmuel show how finfluencers contribute to the financialisation of everyday life by weaving investing into lifestyle narratives, normalising speculative behaviours \cite{HayesFinfluence2024}. This work highlights diverse motivations, ranging from education and community building to self-promotion, and identifies conflicts of interest when creators profit from endorsed products.

Quantitative studies demonstrate that finfluencer activity shapes crowd sentiment. Haase et al.\ find that finfluencers' sentiment Granger-causes broader crowd sentiment, particularly in crypto discussions \cite{HaaseFinfluencersSentiment2025}. Meyer et al.\ document emotional contagion in YouTube crypto content, with influencers' tone mirrored in audience comments \cite{MeyerHighOnBitcoin2023}. These findings suggest finfluencers modulate affect, amplifying herd dynamics in volatile markets.

\subsubsection{Credibility in Finfluencer Research}

Existing research on finfluencer credibility has predominantly adopted a \emph{credential-based} perspective, conceptualising trustworthiness in terms of static attributes such as formal education, professional background, financial certifications, and practitioner experience~\cite{NurhandayaniWhoDeserves2025}. This framing treats credibility as a property that individuals possess by virtue of their qualifications, implicitly privileging institutional markers over situated practices. However, such an approach may not fully capture how credibility operates in decentralised, pseudonymous environments like crypto, where formal credentials are often absent, unverifiable, or irrelevant to community norms. In these contexts, credibility may instead emerge from observable behaviours (such as transparent disclosure, restraint in promotion, and responsiveness to community feedback) rather than from static signals of expertise. This gap motivates our inquiry into how crypto KOLs \emph{enact} credibility through self-determined practices, an approach we analyze empirically in Section~\ref{sec:rq3} and theorise in Section~\ref{sec:discussion}.

\subsubsection{Creator Labour Under Financialisation}

Duffy's work on \emph{aspirational labour} shows how creators invest substantial unpaid time building visibility in hope of future opportunities \cite{DuffyAspirationalWork2017}. Financialisation intensifies these dynamics. Alacovska and Chalcraft analyse NFT artists' work as ``speculative labour,'' where creative outputs become financial assets subject to extreme volatility \cite{AlacovskaSpeculativeLabour2024}. Internet celebrity dynamics \cite{AbidinInternetCelebrity2018} illustrate how rapid visibility generates both opportunities and vulnerability to backlash when advice fails. In such environments, creators' incomes, reputations, and portfolios become entangled with promoted products.

This literature establishes that influencers are not merely communicators but \emph{workers} under conditions of nested precarities. Crypto KOLs inherit these dynamics but operate in a high-risk financial domain where attention fluctuations translate directly into gains or losses, foregrounding questions of credibility and responsibility.

\subsection{Crypto as a Distinct Context}

Cryptocurrency presents a qualitatively different context for influencer work. We synthesise research on usability, regulation, scams, and community trust to explain why crypto KOLs occupy a particularly consequential position.

\subsubsection{Fragile Mental Models and Usability Struggles}

HCI research underscores that cryptocurrencies remain difficult to use safely. Many users struggle with private-key management and rely on third-party services, exposing them to custodial risk \cite{MoserBitcoinUsability2013,Krombholz2016}. Studies reveal fragile mental models of custody, anonymity, and transaction finality, including assumptions about reversibility that rarely hold \cite{Mai2020,Abramova2021}. This suggests that KOL audiences often lack foundational knowledge to critically evaluate advice.

\subsubsection{Regulatory Gaps and Weak Consumer Protections}

Cryptocurrency promotion falls between or outside conventional investor-protection regimes. Finfluencers' undisclosed sponsorships and pump-and-dump schemes evade rules designed for licensed advisors \cite{StefanouFinfluencersReg2022}. Regulatory agencies have issued guidelines targeting misleading promotions \cite{SECInvestorCrypto,FCA2023Finfluencer,ASAInfluencer2023}, yet enforcement remains patchy and jurisdictionally fragmented. Crypto KOLs often operate in even murkier environments, promoting speculative assets that cross national boundaries and challenge securities law. This intensifies both KOLs' impact and ethical tensions around disclosure.

\subsubsection{Scams and Community Trust}

Oak and Shafiq analyse scam-related discourse on Reddit, identifying roles of victims, vigilantes, and advice-givers that underscore grassroots demand for trustworthy guidance \cite{OakScamReddit2025}. On-chain analyses reveal ecosystems of scam tokens and rug pulls that exploit naive investors \cite{Cernera2023rugpulls}. Research on user trust highlights a ``centralised trust in decentralised systems'' paradox: users anchor trust in exchanges and charismatic community figures, precisely the intermediaries decentralisation was meant to obviate \cite{KhairuddinBitcoinMotivations2016,SasKhairuddin2017,BappyCentralizedTrust2025}.

Together, these strands establish that crypto is a socio-technical environment where users struggle with operations, protections are weak, and scams are pervasive. Users often ``outsource'' due diligence to community figures, including KOLs whose endorsements function as trust shortcuts, akin to, but less regulated than, traditional financial advisors.

\subsection{Crypto KOLs as Digital Workers}

We now reframe crypto KOLs as a form of digital creator labour at the intersection of influencer culture and crypto's specific risks.

\subsubsection{Double Exposure: Platform and Market Precarity}

Crypto KOLs exemplify intersecting forms of precarity. Like other creators, they depend on platform algorithms and constant content production \cite{DuffyPlatformPractices2019,CotterVisibilityGame2019}, engaging in visibility and relational labour \cite{AbidinVisibility2016,BaymPlaying2018} while navigating nested precarities \cite{DuffyNestedPrecarities2021}.

However, unlike lifestyle influencers, crypto KOLs' success is tied directly to volatile asset prices. Their work constitutes speculative labour \cite{AlacovskaSpeculativeLabour2024}, with reputations and portfolios entangled with promoted products. Lacking institutional safeguards or licensing, they face a \emph{double exposure} to both platform dynamics and financial markets, rendering their labour uniquely precarious. Internet celebrity dynamics \cite{AbidinInternetCelebrity2018} apply forcefully: rapid visibility generates opportunities but also vulnerability to backlash when projects fail.

\subsubsection{Self-Determination Theory and KOL Motivations}

SDT provides a framework for understanding motivations beyond ``greed'' versus ``altruism.'' Central to SDT is a continuum of relative autonomy, ranging from amotivation through external regulation to fully internalised, intrinsic motivation; the degree to which external incentives become self-endorsed determines the quality and persistence of engagement \cite{RyanDeci2000}. SDT also posits three basic psychological needs: \emph{autonomy}, \emph{competence}, and \emph{relatedness}, whose satisfaction supports this internalisation process and fosters well-being \cite{DeciRyan1985,RyanDeci2000}.

SDT has been applied to online communities and digital work. Nov finds Wikipedians driven by intrinsic enjoyment, learning, and prosocial motives \cite{NovMotivatesWikipedians2007}. Tyack and Mekler show that supporting SDT needs fosters sustained engagement \cite{TyackMeklerSDT2020}. Studies of content creators find that YouTube producers value autonomy, competence, and relatedness, with monetary rewards often secondary \cite{LiuAnxietyCommunities2021,torhonen2019}.

Applying SDT to crypto KOLs, we conceptualise how they may experience \emph{autonomy} in choosing projects, \emph{competence} in decoding protocols for lay audiences, and \emph{relatedness} in cultivating communities. Simultaneously, sponsorships and portfolio performance exert external pressures. An SDT lens distinguishes internalised motivations (education, community stewardship) from controlled ones (sponsor obligations, trading gains), providing a framework for analysing credibility practices.

\subsubsection{Market-Level Influence}

Research links online influence to financial outcomes. Studies of r/WallStreetBets show how social dynamics moved markets during the GameStop squeeze \cite{LucchiniGamestopWSB2022,warkulat2024}. In crypto, Musk's tweets produce significant price changes \cite{AnteMuskCrypto2023}, and finfluencer endorsements channel capital into tokens \cite{Merkley2024,HayesFinfluence2024,HaaseFinfluencersSentiment2025,MeyerHighOnBitcoin2023}.

Across this literature, influencers are modelled as sources of signals correlated with market movements. Much less attention is paid to their experiences as workers, or how motivations and ethical commitments shape their practices.

\subsubsection{Conceptual Summary: Creator-Labour and SDT Concepts Applied to Crypto KOLs}

Table~\ref{tab:concepts-kols} summarises how core concepts from creator-labour scholarship and SDT map onto the work of crypto KOLs. This conceptual mapping guides our empirical analysis by identifying the forms of labour, precarity, and motivation that characterise KOL work.

\begin{table*}[htbp]
\centering
\caption{Conceptual framework mapping creator-labour constructs and SDT needs to the work of crypto key opinion leaders.}
\label{tab:concepts-kols}
\small
\renewcommand{\arraystretch}{1.2}
\begin{tabular}{@{} >{\raggedright\arraybackslash}p{3.5cm} >{\raggedright\arraybackslash}p{6.2cm} >{\raggedright\arraybackslash}p{6.2cm} @{}}
\toprule
\textbf{Concept} & \textbf{Definition} & \textbf{Application to Crypto KOLs} \\
\midrule
\textit{Creator-labour constructs} & & \\[-4pt]
\cmidrule(r){1-1}
\textbf{Visibility labour} {\cite{AbidinVisibility2016}} & Ongoing work of posting, curation, and interaction to maintain public presence & Continuous content production, audience engagement, and cross-platform presence management \\[6pt]
\textbf{Relational labour} {\cite{BaymPlaying2018}} & Effort invested in building and maintaining intimate audience relationships & Cultivating parasocial bonds through community engagement and sharing personal experiences \\[6pt]
\textbf{Aspirational labour} {\cite{DuffyAspirationalWork2017}} & Unpaid or under-compensated work invested in hope of future opportunities & Heavy early-career investment in content and community before monetisation \\[6pt]
\textbf{Nested precarities} {\cite{DuffyNestedPrecarities2021}} & Overlapping insecurities at platform, industry, and economic levels & Algorithmic volatility, market crashes, sponsorship instability, and regulatory uncertainty \\[6pt]
\textbf{Speculative labour} {\cite{AlacovskaSpeculativeLabour2024}} & Creative work entangled with volatile financial assets & Reputations and incomes tied to token prices; portfolios often include promoted projects \\[6pt]
\textbf{Internet celebrity} {\cite{AbidinInternetCelebrity2018}} & Rapid visibility generating opportunities and vulnerabilities & Successful calls attract followers and sponsorships; failed calls invite backlash \\[3pt]
\addlinespace[4pt]
\textit{SDT basic needs}$^{\dagger}$ & & \\[-4pt]
\cmidrule(r){1-1}
\textbf{Autonomy} & Need to feel self-directed and volitional & Freedom to choose projects, content style, and community norms \\[6pt]
\textbf{Competence} & Need to feel effective and capable & Satisfaction from accurate analyses, successful predictions, and audience learning \\[6pt]
\textbf{Relatedness} & Need for meaningful social connections & Community ties and peer networks; prosocial motives coexist with commercial ones \\
\bottomrule
\addlinespace[2pt]
\multicolumn{3}{@{}l}{\footnotesize$^{\dagger}$\,All three needs drawn from the self-determination theory~\cite{DeciRyan1985,RyanDeci2000}.}
\end{tabular}
\end{table*}

\subsection{Research Gap}

A comprehensive mapping of the literatures, their relevance to crypto KOLs, and remaining gaps is provided in \Cref{tab:lit-synthesis} (Appendix). Despite this rich body of work, a crucial gap remains. Existing research focuses on \emph{crypto users and retail investors}, examining their struggles and trust practices; on \emph{market-level consequences}, modelling influencers as sentiment signals; or on \emph{lifestyle influencers and generic finfluencers}, without attending to crypto's specific conditions. We lack an in-depth, SDT-informed qualitative account of \emph{crypto KOLs themselves}.

This gap matters because crypto KOLs occupy a critical, under-researched position as community figures whose motivations and practices have material consequences for followers' financial well-being. Understanding KOLs as workers, rather than merely as market signals, is essential for theory, intervention design, and policy. Our research questions (Section~\ref{sec:introduction}) address this gap: KOL motivations (RQ1), distinguishing practices (RQ2), and credibility and ethical reasoning (RQ3).


\begin{comment}

\begin{figure}
  \centering
  \includestandalone{figures/milestone-timeline}
  \caption{Timeline of key milestones spanning parasocial theory, SDT, and crypto influencer evidence.}
  \label{fig:milestone-timeline}
\end{figure}

\begin{figure}
  \centering
  \includestandalone{figures/state-of-art}
  \caption{State of the art: linking influencer mechanisms, the crypto KOL context, and SDT, highlighting the research gap.}
  \label{fig:state-of-art}
\end{figure}

\end{comment}